\section{Metodología de diseño}

\subsection{Carga axial de servicio}
La carga axial que actúa sobre el suelo incorpora, además de la solicitación transmitida por la superestructura, el peso propio de la zapata y el peso del relleno de suelo dispuesto sobre ella hasta el nivel de sello. Para una zapata de dimensiones $B \times L$ y altura $H$, con sello a profundidad $D_f$, este peso adicional es
\begin{equation}
W = \left( \gamma_s \, D_f + \gamma_h \, H \right) B \, L,
\end{equation}
de modo que la carga axial de servicio sobre el suelo resulta
\begin{equation}
N = |P| + W,
\end{equation}
donde $P$ es la carga axial transmitida por el muro o columna en la combinación correspondiente.

\subsection{Presión de contacto en muros}
En los muros, la presión de contacto se evalúa en una dirección. A partir del momento $M$ y la carga axial $N$ se determina la excentricidad
\begin{equation}
e = \frac{|M|}{N},
\end{equation}
que se compara con el límite del núcleo central $L/6$. Cuando la resultante cae dentro del núcleo ($e \le L/6$), toda la base permanece comprimida y la presión máxima se obtiene mediante la distribución trapezoidal
\begin{equation}
\sigma = \frac{N}{B\,L}\left( 1 + \frac{6e}{L} \right).
\end{equation}
Cuando la excentricidad supera el núcleo ($e > L/6$), parte de la base se despega y la distribución se vuelve triangular, con una longitud comprimida
\begin{equation}
L_p = 3\left( \frac{L}{2} - e \right), \qquad \sigma = \frac{2N}{B\,L_p}.
\end{equation}
En este caso se exige que la zona comprimida abarque al menos el 80\% de la longitud de la zapata, esto es $L_p \ge 0,80\,L$. En ambos casos la presión máxima debe cumplir $\sigma \le \sigma_{adm}$, adoptando el valor admisible estático o sísmico según la combinación.

\subsection{Presión de contacto en muros perimetrales: zapata excéntrica}
\label{subsec:zapata_excentrica}
Los muros perimetrales se sitúan en el límite del edificio, de modo que su zapata no puede proyectarse hacia el exterior de la línea de cierro. En consecuencia, el muro queda en el borde exterior de la zapata y la base se desarrolla únicamente hacia el interior. El centro de la zapata queda así desplazado respecto del eje del muro, introduciendo en la dirección transversal $B$ una excentricidad geométrica
\begin{equation}
e_B = \frac{B}{2} - \frac{e_w}{2} = \frac{B - e_w}{2},
\end{equation}
donde $e_w$ es el espesor del muro. Esta excentricidad geométrica se combina con la excentricidad de carga $e_L = |M|/N$, asociada al momento en el plano del muro, evaluándose la presión de contacto en la esquina más cargada como
\begin{equation}
\sigma_{\max} = \frac{N}{B\,L}\left( 1 + \frac{6 e_B}{B} + \frac{6 e_L}{L} \right),
\end{equation}
verificándose $\sigma_{\max} \le \sigma_{adm}$. Dado que $e_B$ es inherente a la condición de borde y no proviene de la solicitación, en estos elementos el criterio del 80\% de compresión no resulta aplicable; en su lugar se reporta la fracción de base comprimida resultante de la geometría adoptada. El ancho $B$ se ajusta hasta satisfacer simultáneamente la presión admisible y la estabilidad al deslizamiento, manteniendo $L$ igual al largo del muro.

\subsection{Presión de contacto en columnas}
En las columnas, la verificación es biaxial. A partir de los momentos en ambos ejes se determinan las excentricidades
\begin{equation}
e_x = \frac{|M_x|}{N}, \qquad e_y = \frac{|M_y|}{N}.
\end{equation}
Si la resultante se mantiene dentro del núcleo central, condición expresada por $\dfrac{e_x}{B} + \dfrac{e_y}{L} \le \dfrac{1}{6}$, la presión máxima en la esquina más cargada es
\begin{equation}
\sigma = \frac{N}{B\,L}\left( 1 + \frac{6 e_x}{B} + \frac{6 e_y}{L} \right).
\end{equation}
Fuera del núcleo, la presión se obtiene mediante redistribución por cada eje, adoptando el valor máximo resultante. La verificación exige $\sigma \le \sigma_{adm}$.

\subsection{Verificación al corte unidireccional}
La resistencia al corte aportada por el hormigón se evalúa según ACI 318-19 mediante
\begin{equation}
v_c = 0,53\sqrt{f'_c} = 10,01~\mathrm{kg/cm^2}.
\end{equation}
La sección crítica se ubica a una distancia $d$ de la cara del elemento. Con la presión neta mayorada $q_u$ actuando sobre la proyección en voladizo $L_1$, el esfuerzo de corte solicitante y la resistencia de diseño, por unidad de ancho $b_w$, son
\begin{equation}
V_u = q_u\,(L_1 - d)\,b_w, \qquad \phi V_c = \phi \, v_c \, b_w \, d,
\end{equation}
con $\phi = 0,75$. La verificación se satisface cuando $V_u \le \phi V_c$.

\subsection{Verificación al punzonamiento}
En las zapatas aisladas de columna se verifica adicionalmente el punzonamiento en dos direcciones, sobre un perímetro crítico ubicado a $d/2$ de las caras del soporte:
\begin{equation}
b_o = 2\left( b_x + d \right) + 2\left( b_y + d \right).
\end{equation}
La resistencia del hormigón al corte por punzonamiento es
\begin{equation}
v_{c,2D} = 1,06\sqrt{f'_c} = 20,03~\mathrm{kg/cm^2}.
\end{equation}
El esfuerzo solicitante se obtiene descontando la presión que actúa dentro del perímetro crítico, y la resistencia de diseño se evalúa sobre dicho perímetro:
\begin{equation}
V_u = P_u - q_u\,A_{crit}, \qquad \phi V_c = \phi \, v_{c,2D} \, b_o \, d,
\end{equation}
verificándose $V_u \le \phi V_c$.

\subsection{Estabilidad al deslizamiento}
La estabilidad frente al deslizamiento se evalúa mediante el factor de seguridad
\begin{equation}
FS_d = \frac{\mu \, N}{Q} \ge 1,5,
\end{equation}
donde $\mu = 0,40$ es el coeficiente de roce suelo-fundación y $Q$ es la fuerza horizontal solicitante. En los muros, $Q$ corresponde al corte en el plano, mientras que en las columnas se compone vectorialmente a partir de las componentes en ambos ejes:
\begin{equation}
Q = |V| \quad \text{(muros)}, \qquad Q = \sqrt{V_2^2 + V_3^2} \quad \text{(columnas)}.
\end{equation}
En los muros perimetrales con zapata excéntrica, esta verificación resulta determinante en el ajuste del ancho $B$, dado que la holgura disponible en presión de contacto es mayor que en deslizamiento.

\subsection{Diseño de la armadura de flexión}
El momento de diseño se obtiene tratando la zapata como un voladizo empotrado en la cara del elemento, con una longitud de volado
\begin{equation}
L_v = \frac{L - b}{2},
\end{equation}
donde $b$ es la dimensión del muro o columna en la dirección considerada. El momento solicitante por unidad de ancho es
\begin{equation}
M_u = \frac{q_u \, L_v^2}{2}.
\end{equation}
En los muros perimetrales con zapata excéntrica, el voladizo se desarrolla en un único sentido hacia el interior, con $L_v = B - e_w$. El área de acero requerida se determina iterativamente mediante el bloque de compresión equivalente de Whitney:
\begin{equation}
A_s = \frac{M_u}{\phi \, f_y \left( d - \dfrac{a}{2} \right)}, \qquad a = \frac{A_s \, f_y}{0,85 \, f'_c \, b},
\end{equation}
con $\phi = 0,90$. El proceso parte de un valor inicial de $a$ y se repite hasta la convergencia del par $(A_s, a)$. El área resultante se compara con la armadura mínima $A_{s,\min} = 21,6~\mathrm{cm^2/m}$, adoptándose el mayor de ambos valores. La necesidad de disponer una parrilla de refuerzo se establece comparando la tensión de tracción resultante en la base con un valor límite de $8,5~\mathrm{kg/cm^2}$: cuando se supera dicho límite, el elemento requiere una parrilla diseñada por flexión; en caso contrario, gobierna la armadura mínima. A partir del área adoptada se selecciona el diámetro y la separación de las barras.